% ---------------------------------------------------------------------
% 第 1 章 引言
% ---------------------------------------------------------------------
\section{引言}

\subsection{背景与动机}

刚体位姿（姿态 + 位置）的参数化是机器人控制的起点。单位对偶四元数将两者装入一个 8 维代数对象，运算全局无奇异，且乘法直接实现位姿复合；\cite{P2} 在此参数化上给出了带 $L_2$ 扰动衰减保证的运动学跟踪控制器，是 DQ 控制的代表性结果。Cohen \& Shoham \cite{P1} 进一步引入超对偶四元数（HDQ）：给 DQ 增加一个幂零单位 $\varepsilon^*$（$\varepsilon^{*2}=0$），使 HDQ 能同时表示位姿与其一阶时间导数，链式乘法自动微分，从而正运动学一次连乘同时输出位姿与 twist。

但 HDQ 到动力学层面仍差一阶：计算力矩控制需要任务空间加速度 $\dot{\boldsymbol\xi}$ 与雅可比导数项 $\dot J\dot{\boldsymbol q}$，而 \cite{P1} 的 HDQ 结构中两个幂零单位（$\varepsilon$ 承载平移、$\varepsilon^*$ 承载一阶导）均已占用，$\varepsilon\varepsilon^*$ 项只表示“平移分量的一阶导数”，不含二阶时间导数。更重要的是\textbf{误差体系}：现有 DQ 框架的误差对象只有位姿一层，扰动模型假设速度级加性扰动且属于 $L_2$；进入动力学接口后出现三个结构性缺口，即位姿与速度误差分离（\cite{Ch20} 类律的位姿反馈取螺旋对数，其导数映射在大误差处奇异）、加速度层扰动无入口、偏差型不确定性不满足 $L_2$ 假设，既有证书随之失效。

\subsection{本文贡献}

\begin{enumerate}[leftmargin=2em]
  \item \textbf{TODQ 运动学}：基于三项多项式代数构造三阶对偶四元数串联运动学，突破现有 HDQ 正运动学止于一阶导数的限制，一次 $O(n)$ 连乘同时输出位姿、twist 与任务空间加速度，为加速度级控制提供解析前馈。
  \item \textbf{一体化误差体系}：在两项 HDQ 元素上定义误差对象，一次乘法同时生成右不变位姿误差与几何一致 twist 误差，并证明闭环误差动力学为严格级联标准形。
  \item \textbf{控制律与精确证书}：提出 $A^\top$ 整形计算力矩律，闭环存储函数满足精确耗散等式，统一给出无扰收敛、$L_2$ 衰减充要判据与偏差型扰动的均方界；$L_2$ 增益经最优化恰为最小阻尼的倒数，与线性化真实 H∞ 范数相等，衰减指标可直接换算为阻尼下界。
  \item \textbf{实机验证}：在六自由度机械臂上完成四律同预算对比与 $\gamma$ 增益阶梯共 33 次运行，各证书全程成立，并区分任务律控制的瞬态精度与拴定弹簧、摩擦主导的准静态残差。
\end{enumerate}

TODQ 是 \cite{P1} HDQ（幂零单位承载导数）向二阶的最小扩展 \cite{Con25}；0 阶误差退化为 \cite{P2}（§4.4）。
